\documentclass[../../../include/open-logic-section]{subfiles}

\begin{document}
	
\olfileid{sth}{z}{nat}
\olsection{మన సహజ సంఖ్యల ఎంపిక}

%అనౌపచారిక \stagesinf{} సూత్రాన్ని
%సమర్థించవచ్చని అనుకుందాం. అయితే దాని నుంచి
%పొందిన ప్రత్యేక స్వీకృతంపై కూడా వ్యాఖ్యానించడం విలువైనది.

\olref[infinity-again]{defnomega}లో
``సహజ సంఖ్యలు'' అనే
వ్యక్తీకరణను స్పష్టంగా
నిర్వచించాం. ఈ నిర్దేశాన్ని
మీరు ఎలా అర్థం చేసుకోవాలి?
ఇది తాత్త్వికంగా వస్తువులు
ఏమిటన్న వాదన కాదు;
కొన్ని సమితులను సహజ
సంఖ్యలు\emph{గా పరిగణించాలనే}
నిర్ణయం మాత్రమే. అలా
భావించడానికి గల కారణాలను
\olref[sfr][rel][ref]{sec},
\olref[sfr][arith][ref]{sec},
\olref[sfr][infinite][dedekindsproof]{sec}లో
స్పృశించాం. అయితే ఆ
కారణాలను మరింత స్పష్టంగా
చెప్పవచ్చు.

మన అనంతత్వ స్వీకృతం
\citet{VonNeumann1925}ను
అనుసరిస్తుంది. కానీ
దానికి బదులుగా ఈ
కింది స్వీకృతాన్ని
ఎంచుకోవచ్చు:

\begin{defish}
\emph{సెర్మెలో \citeyear{Zermelo1908Untersuchungen}
అనంతత్వ స్వీకృతం.}
$\emptyset \in A$ ఉండి,
$(\forall x \in A)\{x\} \in A$
అయ్యే ఒక సమితి $A$
ఉంటుంది.
\end{defish}

మన (వాన్ న్యూమన్ స్ఫూర్తిగల)
అనంతత్వ స్వీకృతానికి
బదులు సెర్మెలో స్వీకృతాన్ని
వాడినా, డెడెకిండ్ అనంత
సమితి, దానితో డెడెకిండ్
బీజగణితం లభించేవి.
సెర్మెలో పద్ధతిలో మన
బీజగణితపు ప్రత్యేక మూలకం
మళ్లీ $\emptyset$గానే
(మన $0$కు ప్రతినిధిగా)
ఉండేది. కానీ ఒకటి-ఒకటి
ప్రమేయం $x \mapsto x \cup \{x\}$
బదులు $x \mapsto \{x\}$
అనే ప్రమేయంతో ఇవ్వబడేది.
దీని సరళమైన ఫలితం:
సెర్మెలో $2$ను
$\{\{\emptyset\}\}$గా పరిగణిస్తాడు;
మనమైతే (వాన్ న్యూమన్‌తో)
$2$ను $\{\emptyset, \{\emptyset\}\}$గా
పరిగణిస్తాం.

ఒక అనంతత్వ స్వీకృతం
బదులు మరొకదాన్ని ఎందుకు
ఎంచుకోవాలి? ముఖ్యమైన
ఆచరణాత్మక కారణం:
వాన్ న్యూమన్ పద్ధతి
అతిపరిమిత సంఖ్యలను
చక్కగా నిర్వహించేలా
``విస్తరిస్తుంది''.
దీనిని \olref[ordinals][]{chap}
నుంచి పరిశీలిస్తాం.
అయితే \emph{అంకగణితం
చేయడం} అనే సాధారణ
దృష్టిలో రెండు పద్ధతులూ
సమానంగా పనికొస్తాయి.
కాబట్టి సహజ సంఖ్యలు
సమితులే \emph{అని}
ఎవరైనా చెబితే,
సహజమైన ప్రశ్న:
\emph{అవి ఏ సమితులు?}

ఈ ఖచ్చితమైన ప్రశ్నను
\citet{Benacerraf1965}
ప్రసిద్ధం చేశాడు. అయితే
మనం ఇప్పటికే పలుమార్లు
ఎదుర్కొన్న ఒక విషయానికి
ఇది ప్రసిద్ధ ఉదాహరణ
మాత్రమే అని గుర్తుంచుకోవాలి.
మూల విషయం ఇది:
వాస్తవ సంఖ్యలు, కరణీయ
సంఖ్యలు, పూర్ణ సంఖ్యలు,
సహజ సంఖ్యలు వంటి
``సహజంగా భావించే''
వస్తు రకాలనూ, అలాగే
క్రమయుగ్మాలు, ప్రమేయాలు,
సంబంధాలనూ \emph{అనుకరించే}
మార్గాన్ని సమితి సిద్ధాంతం
ఇస్తుంది. అయితే ఆ
అనుకరణకు సమితి సిద్ధాంతం
ఎప్పుడూ \emph{ఏకైక}
ఎంపికను ఇవ్వదు. సమానంగా
ఉపయోగపడే ప్రత్యామ్నాయాలు
\emph{ఎల్లప్పుడూ} ఉంటాయి.

\end{document}
